\scanpage{53}


\margin{$\operatorname{Hom}(x,Sy)=\operatorname{Hom}(Tx,y)$}

% letterhead
\begin{center}
INSTITUT\\
DES HAUTES ÉTUDES SCIENTIFIQUES\\
5, ROND-POINT BUGEAUD, PARIS XVI\textsuperscript{e}\\
KLÉBER 48-56
\end{center}

\uline{Prop}

\diagram{\begin{tikzcd}[ampersand replacement=\&,column sep=large]
A\arrow[r,shift left,"T"] \& B\arrow[l,shift left,dashed,"S"]
\end{tikzcd}}

$T$ foncteur ``passage au quotient'' \unsure{par|pour} une catégorie de fractions.

Soit $x\in\mathrm{Ob}\,A$, $y=T(x)$. Propriétés équivalentes

a) $x$ est libre au dessus de $y$, i.e. $\forall\,x'\in\mathrm{Ob}\,A$,
\[
\operatorname{Hom}_{A}(x',x)\xrightarrow{\sim}\operatorname{Hom}_{B}(Tx',Tx).
\]

b) Si $x'\unsure{\to|-}x''$ dans $A$ est un ``$\deleted{T}C$ isomorphisme'', (i.e. $Tx'\unsure{-|\to}Tx''$ un isomorphisme) alors
\[
\operatorname{Hom}_{A}(x'',x)\unsure{\to|-}\operatorname{Hom}_{A}(x',x)
\]
est une bijection.

c) [Si $B=A/C$, $A$ catégorie abélienne, $C$ sous-catégorie abélienne épaisse] Tout sous-objet de $x$ appartenant à $\mathrm{Ob}\,C$ est \uline{nul}, et $x$ est facteur direct dans tt $y\supset x$ tel que $y/x\in\mathrm{Ob}\,C$

d) [Si l'adjoint $S$ de $T$ existe]
\[
x\xrightarrow{\unsure{\mathrm{can}|\infty|\sim}}STx\quad\text{est un isomorphisme}
\]

e) \inserted{[Si l'adjoint $S$ de $T$ existe]} $x$ est dans l'image essentielle de $S$, i.e. isomorphe à un objet de la forme $Tz\ [z\in\mathrm{Ob}\,C]$

% archive sheet number: 52
